\chapter{Первый взлёт}\label{ch:vzlet}

\textbf{Что вы узнете:} как выглядит самое короткое шоу, что такое шаг, откуда берутся секунды.

\section{Самое короткое шоу}

\solfile{l02_first.py}

Разберём построчно.
\begin{itemize}
  \item \kod{from litebee_show import Show} берёт из библиотеки слово \kod{Show}. Эта строка всегда первая.
  \item \kod{show = Show('Первый взлёт')} создаёт шоу и даёт ему имя. Переменная \emph{обязана} называться \kod{show}: приложение берёт шоу именно из неё.
  \item \kod{d = show.drone(1)} создаёт дрона с номером 1. Дальше \kod{d} — это «наш дрон»: команды вида \kod{d.что-то(...)} добавляют ему шаги.
  \item \kod{d.calibrate(5)}: калибровка, 5 секунд. Обычно это первый шаг.
  \item \kod{d.takeoff(100, 5)}: взлёт на высоту 100\cm за 5 секунд.
  \item \kod{d.land()}: посадка.
\end{itemize}

Строки, которые начинаются с \texttt{\#}, — комментарии. Python их пропускает, они для людей.

\section{Шаги и секунды}

Каждая команда добавляет дрону один \textbf{шаг}. Шаги выполняются строго по порядку: следующий начинается, когда закончился предыдущий. Число секунд в команде — это \textbf{длительность шага}, а не момент времени. \kod{d.takeoff(100, 5)} значит «взлетай пять секунд», а не «будь наверху на пятой секунде».

\begin{figure}[H]
  \centering
  % Лента l02_first.py: 5 + 5 + 3 = 13 секунд
\begin{tikzpicture}[x=0.9cm, font=\sffamily\small]
  \def\h{0.9}
  \foreach \a/\b/\t/\f in {0/5/{calibrate(5)}/paper, 5/10/{takeoff(100, 5)}/accentsoft, 10/13/{land()}/paper}{
    \draw[fill=\f, draw=ink!60, thick] (\a,0) rectangle (\b,\h);
    \node[font=\scriptsize, text=ink] at ({(\a+\b)/2},{\h/2}) {\t};
  }
  \draw[-{Stealth}, ink] (0,-0.5) -- (13.9,-0.5) node[right] {$t$, с};
  \foreach \t in {0,5,10,13}{\draw[ink] (\t,-0.42) -- (\t,-0.58); \node[below, font=\scriptsize] at (\t,-0.58) {\t};}
\end{tikzpicture}

  \caption{Лента времени шоу \texttt{l02\_first.py}: $5 + 5 + 3 = 13$ секунд.}\label{fig:tl-first}
\end{figure}

Время программы дрона — сумма времён шагов. Посадка в симуляторе всегда занимает 3 секунды, хотя в скобках у \kod{land()} ничего не написано.

\begin{zametka}
У каждого дрона свои часы, они начинаются с нуля. Дрон не знает, что делают другие (урок~\ref{ch:ochered}).
\end{zametka}

\section{Лимит высоты}

Высота — расстояние от пола до дрона в сантиметрах. Выше лимита лететь нельзя: по умолчанию лимит 150\cm. Его можно поменять на верхней панели, \knopka{Лимит высоты}, от 20 до 200\cm. Посмотрите, что скажет проверка, если взлететь слишком высоко:

\solfile{bad_examples/bad_height.py}

Нажмите \knopka{Проверить}: в списке появится ошибка «Дрон 1: высота 180 см больше лимита 150 см, полёт будет отклонён».

\begin{probuy}
  \item Взлетите на 60\cm за 4 секунды и сразу сядьте.
  \item Добавьте зелёный свет (\kod{d.rgb('green')}), чтобы дрон повисел 3 секунды перед посадкой. Подсказка: ожидание записывается как \kod{d.set_z(100, 3)} (о нём урок~\ref{ch:vysota}); потом верните светодиод в \kod{'black'}.
  \item Сколько секунд длится шоу из файла \texttt{l02\_first.py}, если взлёт записать как \kod{takeoff(100, 8)}?
  \item Что нужно поменять в файле \texttt{bad\_height.py}, чтобы проверка прошла?
\end{probuy}
\otvety{2}
